% !TeX program = xelatex
\documentclass[11pt, a4paper]{article}

\usepackage[a4paper, top=2.5cm, bottom=2.5cm, left=2.8cm, right=2.2cm]{geometry}
\usepackage{fontspec}
\usepackage[vietnamese, english, bidi=basic, provide=*]{babel}
\babelprovide[import, onchar=ids fonts]{english}

% Font chữ chuẩn cho Overleaf (XeLaTeX / LuaLaTeX)
\IfFontExistsTF{Times New Roman}{\babelfont{rm}{Times New Roman}}{\babelfont{rm}{TeX Gyre Termes}}

\usepackage{enumitem}
\setlist[itemize]{label=-}
\usepackage{amsmath}
\usepackage{amssymb}
\usepackage{booktabs}
\usepackage{array}
\usepackage{longtable}
\usepackage{tabularx}
\usepackage{float}
\usepackage{xcolor}
\usepackage{graphicx}
\usepackage{hyperref}

% Định dạng màu sắc chuẩn học thuật: chữ đen tuyệt đối
\definecolor{textblack}{HTML}{000000}
\color{textblack}
\hypersetup{
    colorlinks=true,
    linkcolor=black,
    citecolor=black,
    urlcolor=black,
    pdftitle={Báo cáo cuối kỳ Trí tuệ nhân tạo nâng cao - Game AI Connect 4},
    pdfauthor={Nguyễn Hoàng Khang}
}

\linespread{1.25}
\setlength{\parindent}{0pt}
\setlength{\parskip}{6pt}

\begin{document}

% ==============================================================================
% TRANG BÌA (TITLE PAGE)
% ==============================================================================
\begin{titlepage}
    \centering
    \vspace*{0.5cm}
    {\large \textbf{[CẦN BỔ SUNG: TÊN TRƯỜNG ĐẠI HỌC / HỌC VIỆN]}}\\[0.2cm]
    {\normalsize \textbf{[CẦN BỔ SUNG: TÊN KHOA / BỘ MÔN]}}\\[0.4cm]
    \rule{0.6\textwidth}{0.6pt}

    \vspace{2.0cm}
    {\large \textbf{BÁO CÁO ĐỒ ÁN CUỐI KỲ}}\\[0.3cm]
    {\large \textbf{HỌC PHẦN: TRÍ TUỆ NHÂN TẠO NÂNG CAO}}

    \vspace{1.5cm}
    \rule{\textwidth}{1.2pt}\\[0.4cm]
    {\Large \textbf{XÂY DỰNG VÀ ĐÁNH GIÁ HỆ THỐNG GAME AI ĐA CẤP ĐỘ CHO TRÒ CHƠI CONNECT FOUR ($7 \times 6$)}}\\[0.4cm]
    \rule{\textwidth}{1.2pt}

    \vspace{2.2cm}
    \begin{tabular}{@{} p{5.2cm} p{8.8cm} @{}}
        \textbf{Sinh viên thực hiện:} & Nguyễn Hoàng Khang \\
        \textbf{Mã số sinh viên:} & 3123410152 \\
        \textbf{Mã nhóm:} & \texttt{[CẦN BỔ SUNG: Mã nhóm]} \\
        \textbf{Giảng viên hướng dẫn:} & \texttt{[CẦN BỔ SUNG: Họ và tên Giảng viên]} \\
        \textbf{Kho lưu trữ mã nguồn:} & \texttt{[CẦN BỔ SUNG: Liên kết GitHub repository]} \\
        \textbf{Mã commit đánh giá cuối:} & \texttt{d90ecfa1e3c9a37367b59f5e010b0de247d09390} \\
        \textbf{Trạng thái thực nghiệm:} & \textbf{Các đợt chạy chính đã báo cáo đều hoàn tất} \\
    \end{tabular}

    \vfill
    {\normalsize Thành phố Hồ Chí Minh, ngày 23 tháng 11 năm 2026}
\end{titlepage}

% ==============================================================================
% TÓM TẮT (ABSTRACT) & MỤC LỤC
% ==============================================================================
\section*{Tóm tắt}
\addcontentsline{toc}{section}{Tóm tắt}

Trong nghiên cứu Trí tuệ nhân tạo (Game AI), việc đánh giá năng lực của các tác tử chơi trò chơi đối kháng thông qua giải đấu nội bộ thường chỉ phản ánh thứ hạng tương đối mà không định lượng được khoảng cách thực tế tới chiến lược tối ưu lý thuyết. Đồ án này xây dựng hệ thống tác tử đa cấp độ cho Connect Four $7 \times 6$, gồm tìm kiếm cổ điển, CNN học có giám sát (CNN v1 và CNN v3 trên 30.000 vị trí có nhãn Pons), policy-gradient và các phương pháp kết hợp mạng với tìm kiếm lấy cảm hứng từ AlphaGo/AlphaZero. Hai run Phase C V2 và round-robin cập nhật đã hoàn tất; một số ablation mở rộng chưa được thực hiện. Warm-start giảm chất lượng validation sau self-play; mức tăng của scratch chưa phân biệt được thống kê.

Quy trình đánh giá kết hợp chặt chẽ giữa thi đấu vòng tròn đối kháng có kiểm soát (28 cặp, 100 ván/cặp, tổng 2.800 ván, ghép cặp thế khai cuộc công bằng, đổi lượt đi trước/sau) và đối chiếu trực tiếp trên tập kiểm thử solver đóng băng gồm 1.500 vị trí độc lập trích xuất từ bộ giải tuyệt đối Pons. Toàn bộ các chỉ số báo cáo đều kèm theo cỡ mẫu và khoảng tin cậy (CI) $95\%$ phân cụm hoặc bootstrap theo khai cuộc ghép cặp; trường hợp khoảng tin cậy chứa giá trị trung hòa (0 đối với sai khác hoặc 0,5 đối với score) được kết luận trung thực là ``chưa phân biệt được''.

Kết quả thực nghiệm chính thức cho thấy:
\begin{enumerate}
    \item \textbf{Nhóm tìm kiếm cổ điển:} Alpha-Beta tăng tiến tỷ lệ nước tối ưu chính xác (\textit{exact}) từ $70{,}33\%$ $[68{,}12;\; 72{,}57]\%$ (độ sâu 1) lên cực đại $83{,}07\%$ $[81{,}18;\; 84{,}96]\%$ (độ sâu 5) và bão hòa ở $81{,}47\%$ $[79{,}46;\; 83{,}37]\%$ (độ sâu 6), vượt trội Heuristic ($73{,}53\%$ $[71{,}34;\; 75{,}66]\%$) và MCTS-500 ($75{,}40\%$ $[73{,}23;\; 77{,}56]\%$, WDL $90{,}80\%$ $[89{,}25;\; 92{,}24]\%$). Random chỉ đạt $37{,}27\%$ $[34{,}83;\; 39{,}72]\%$ exact và $70{,}67\%$ $[68{,}24;\; 73{,}07]\%$ WDL.
    \item \textbf{CNN v1 và CNN v3:} CNN v1 là mốc cơ sở trên dữ liệu nhỏ. CNN v3 dùng 30.000 vị trí; trên validation, Pons-soft đạt exact trung bình $80{,}09\%$ (SD giữa seed $0{,}29$ điểm phần trăm), so với $76{,}93\%$ cho Pons-hard và $73{,}26\%$ cho nhãn Alpha-Beta depth 3. Trên test đóng băng 1.500 vị trí, CNN v3 thô đạt $68{,}60\%$ $[66{,}15;\; 71{,}06]\%$ exact và $87{,}13\%$ WDL; có safety filter đạt $76{,}73\%$ $[74{,}54;\; 78{,}91]\%$ exact và $91{,}00\%$ $[89{,}47;\; 92{,}44]\%$ WDL.
    \item \textbf{PUCT:} Tác tử cấp 7 baseline (CNN v1 + 500 mô phỏng) đạt $78{,}93\%$ $[76{,}80;\; 80{,}93]\%$ exact. Với CNN v3, PUCT-400 warm vòng 0 có filter đạt $82{,}13\%$ $[80{,}14;\; 84{,}03]\%$ exact và $94{,}00\%$ $[92{,}70;\; 95{,}20]\%$ WDL; bản thô đạt $81{,}20\%$ $[79{,}14;\; 83{,}19]\%$ exact.
    \item \textbf{Thực nghiệm tự chơi (Self-play Phase C, hoàn thành 40 vòng lặp):} Cả hai nhánh Warm-start và Scratch đã chạy đủ 40 vòng lặp. Quá trình tự đấu không cải thiện sức chơi ở quy mô thử nghiệm này: nhánh warm-start suy giảm trên validation từ $89{,}53\%$ xuống $85{,}90\%$ (có lọc) và $88{,}46\%$ xuống $78{,}26\%$ (thô); nhánh scratch tăng policy thô từ $33{,}67\%$ lên $47{,}79\%$ nhưng PUCT thô suy giảm từ $81{,}45\%$ xuống $75{,}58\%$, trong khi PUCT có lọc không phân biệt được ưu thế ($81{,}48\%$ lên $82{,}48\%$, CI ghép cặp chứa 0).
    \item \textbf{Tác tử học tăng cường (RL):} Tác tử RL huấn luyện chỉ đối đầu với Random đạt score $0{,}990$ trước Random nhưng sụp đổ trước Heuristic (score $0{,}110$ do Heuristic thắng $86\%$) và đạt score $0{,}000$ trước PUCT.
    \item \textbf{Lợi thế đi trước (RQ6):} Trong 2.800 ván đấu của giải round-robin chính thức, bên đi trước đạt score trung bình $0{,}5325$, khoảng tin cậy bootstrap 95\% [0,5039; 0,5607] (thắng 1.491 ván, chiếm 53,25%; hòa 0; thua 1.219 ván, chiếm 43,54%), xác nhận ưu thế đi trước có ý nghĩa thống kê.
\end{enumerate}

\newpage
\tableofcontents
\newpage

% ==============================================================================
% CHƯƠNG 1. GIỚI THIỆU
% ==============================================================================
\section{Giới thiệu}

\subsection{Bối cảnh nghiên cứu}
Trong lịch sử phát triển của Trí tuệ nhân tạo (AI), các trò chơi đối kháng có thông tin hoàn hảo (perfect-information games) luôn giữ vai trò là môi trường chuẩn mực để thử nghiệm và đánh giá các thuật toán ra quyết định \cite{russell2021artificial}. Việc lựa chọn môi trường thực nghiệm đòi hỏi sự cân nhắc giữa độ phức tạp chiến thuật và khả năng kiểm soát tài nguyên: các trò chơi quá đơn giản (như Tic-Tac-Toe) không đủ độ sâu để bộc lộ hạn chế của thuật toán, trong khi các trò chơi có không gian trạng thái khổng lồ (như Cờ Vua hay Cờ Vây) lại đòi hỏi cụm máy chủ công nghiệp và không có lời giải lý thuyết để kiểm chứng tính tối ưu tuyệt đối.

Connect Four (kích thước chuẩn $7 \times 6$) là môi trường thực nghiệm lý tưởng: trò chơi có khoảng $4{,}5 \times 10^{12}$ trạng thái hợp lệ \cite{tromp2008solving}, đủ phong phú để thách thức các thuật toán tìm kiếm và học sâu, đồng thời đã được chứng minh về mặt toán học là bên đi trước luôn thắng nếu đi hoàn hảo \cite{allis1988knowledge}. Nhờ các bộ giải chính xác mã nguồn mở tốc độ cao như solver của Pascal Pons \cite{pons2019connect4}, mọi quyết định của tác tử đều có thể đối chiếu trực tiếp với chân lý lý thuyết (ground-truth).

\subsection{Trò chơi Connect Four và mô hình hóa bài toán}
Connect Four là trò chơi hai người có tổng bằng không (zero-sum), thông tin hoàn hảo, diễn ra trên bảng đứng kích thước 6 hàng $\times$ 7 cột. Hai người chơi lần lượt thả quân (người 1 mang quân Đỏ, người 2 mang quân Vàng) vào một trong 7 cột chưa đầy; quân cờ rơi xuống vị trí trống thấp nhất của cột do trọng lực. Trò chơi kết thúc khi có một bên tạo thành đường thẳng gồm 4 quân liên tiếp (ngang, dọc hoặc chéo) hoặc bảng cờ đầy (kết quả hòa).

Không gian trạng thái hợp lệ của bài toán được mô hình hóa theo quá trình ra quyết định Markov hữu hạn (finite MDP):
\begin{equation}
\mathcal{M} = \langle \mathcal{S}, \mathcal{A}, \mathcal{P}, \mathcal{R}, \gamma \rangle,
\end{equation}
trong đó $\mathcal{S}$ là tập các cấu hình bàn cờ hợp lệ ($|\mathcal{S}| \approx 4{,}53 \times 10^{12}$); $\mathcal{A} = \{0, 1, 2, 3, 4, 5, 6\}$ là tập 7 cột hành động; $\mathcal{P}(s' \mid s, a)$ là hàm chuyển đổi trạng thái tất định; hàm phần thưởng $\mathcal{R} \in \{+1, -1, 0\}$ chỉ xuất hiện tại trạng thái kết thúc (terminal); và hệ số chiết khấu $\gamma = 1$.

\subsection{Vấn đề nghiên cứu}
Đa số các nghiên cứu thực nghiệm Game AI đối mặt với ba thách thức lớn:
\begin{enumerate}
    \item \textbf{Thước đo tương đối thiếu kiểm chứng khách quan:} Đánh giá sức mạnh tác tử bằng tỷ lệ thắng trong giải đấu nội bộ dễ dẫn đến hiện tượng chu trình (A thắng B, B thắng C, C thắng A) và không cho biết tác tử còn cách chiến lược tối ưu bao xa.
    \item \textbf{Nhầm lẫn giữa năng lực học sâu và quy tắc kinh nghiệm:} Nhiều mô hình học sâu kết hợp quy tắc khẩn cấp (như tự chọn nước thắng ngay hoặc chặn đối thủ) nhưng không bóc tách đóng góp của mạng nơ-ron thuần túy với bộ lọc an toàn, dẫn đến việc đánh giá quá cao khả năng suy luận của mạng.
    \item \textbf{Thiếu kiểm soát sai lệch thực nghiệm:} Nguy cơ rò rỉ dữ liệu giữa ván huấn luyện và ván kiểm thử, hiện tượng lặp lại ván đấu giữa các tác tử tất định, hoặc thiên vị thế cờ khai cuộc chưa được xử lý triệt để trong các báo cáo sinh viên.
\end{enumerate}

\subsection{Mục tiêu và Câu hỏi nghiên cứu}
Đồ án đặt ra 6 câu hỏi nghiên cứu cụ thể nhằm làm sáng tỏ các vấn đề trên:
\begin{description}[leftmargin=1.5cm,style=nextline]
    \item[\textbf{RQ1 (Tương quan giữa sức mạnh thực nghiệm và tính tối ưu):}] Thứ hạng điểm số trong thi đấu đối kháng trực tiếp tương quan như thế nào với tỷ lệ chọn nước tối ưu đo trên bộ giải chuẩn?
    \item[\textbf{RQ2 (Đánh đổi giữa chất lượng quyết định và chi phí tính toán):}] Khi kiểm soát ngân sách suy luận (độ sâu, số mô phỏng, thời gian tính toán), các phương pháp tìm kiếm cổ điển, mạng nơ-ron học có giám sát và tìm kiếm kết hợp mạng đạt hiệu năng ra sao?
    \item[\textbf{RQ3 (Đóng góp của bộ lọc an toàn chiến thuật):}] Quy tắc hậu xử lý (chọn nước thắng tức thì, chặn nước thua tức thì) đóng góp bao nhiêu điểm phần trăm vào tỷ lệ tối ưu của mạng nơ-ron và tìm kiếm cây?
    \item[\textbf{RQ4 (Ảnh hưởng của quy mô dữ liệu và loại nhãn huấn luyện):}] Mở rộng quy mô dữ liệu từ 1.300 lên 30.000 mẫu và thay đổi loại nhãn (Pons-soft, Pons-hard, AB3 teacher) tác động thế nào đến khả năng khái quát hóa của CNN?
    \item[\textbf{RQ5 (Hiệu quả thực nghiệm của vòng lặp tự chơi AlphaZero-inspired):}] Quá trình tự chơi (self-play) kết hợp PUCT có cải thiện được độ chính xác lý thuyết và năng lực chiến đấu so với mô hình tiền huấn luyện có giám sát hay không ở quy mô tài nguyên nhỏ?
    \item[\textbf{RQ6 (Lợi thế của người đi trước):}] Quyền đi trước mang lại ưu thế định lượng cụ thể là bao nhiêu điểm số đối đầu đối với từng họ thuật toán và trên toàn giải đấu?
\end{description}

\subsection{Phạm vi thực hiện và Sai lệch giao thức}
Hệ thống hoàn thiện bao gồm: lõi mô phỏng Connect Four tối ưu hóa bằng Bitboard, bộ kiểm thử tự động, giao diện tương tác người--máy, bộ công cụ trích xuất dữ liệu từ solver, quy trình huấn luyện học sâu và khung đánh giá thống kê. Hai run Phase C V2 và round-robin cập nhật đã hoàn tất. Một số phân tích mở rộng nêu ở phần hạn chế (sweep số mô phỏng, rollout đối chứng, ablation tốc độ học) chưa được chạy; trạng thái hoàn tất ở đây chỉ áp dụng cho các run và đánh giá được báo cáo.

\textbf{Ghi nhận sai lệch giao thức (Protocol Deviation):} Theo kế hoạch ban đầu đăng ký trong \texttt{docs/CRITERIA\_V2.md} và \texttt{docs/COLAB.md}, Phase C dự kiến chạy trên Google Colab GPU (NVIDIA T4). Tuy nhiên, do giới hạn ngắt kết nối phiên làm việc dài (session timeouts) trên môi trường Colab miễn phí, toàn bộ 40 vòng lặp self-play của Phase C đã được chuyển sang thực thi liên tục trên máy trạm cục bộ (NVIDIA GeForce RTX 3050 Laptop GPU, VRAM 4 GB). Mọi đánh giá solver đều thực hiện trên cùng tập test đóng băng 1.500 vị trí và kiểm soát seed ngẫu nhiên.

% ==============================================================================
% CHƯƠNG 2. CƠ SỞ LÝ THUYẾT VÀ CÔNG TRÌNH LIÊN QUAN
% ==============================================================================
\section{Cơ sở lý thuyết và Công trình liên quan}

\subsection{Các thuật toán tìm kiếm đối kháng cổ điển}
Trong lý thuyết trò chơi đối kháng hai người, giải thuật Minimax xác định giá trị tối ưu của trò chơi thông qua đệ quy:
\begin{equation}
V(s) = \begin{cases}
\mathcal{R}(s), & \text{nếu } s \text{ là trạng thái kết thúc}, \\
\max_{a \in \mathcal{A}(s)} V(\mathcal{P}(s, a)), & \text{nếu lượt người chơi } 1, \\
\min_{a \in \mathcal{A}(s)} V(\mathcal{P}(s, a)), & \text{nếu lượt người chơi } 2.
\end{cases}
\end{equation}
Giải thuật Cắt tỉa Alpha-Beta (Alpha-Beta Pruning) tối ưu hóa Minimax bằng cách duy trì hai ngưỡng $[\alpha, \beta]$ biểu thị giá trị tối thiểu người chơi Max đã đảm bảo và giá trị tối đa người chơi Min cho phép; bất kỳ nhánh con nào có giá trị vi phạm khoảng này đều bị cắt bỏ, giảm độ phức tạp từ $\mathcal{O}(b^d)$ xuống $\mathcal{O}(b^{d/2})$ trong trường hợp sắp xếp nước đi tối ưu \cite{russell2021artificial}.

Thuật toán Tìm kiếm cây Monte Carlo (MCTS) dựa trên công thức UCT (Upper Confidence Bound for Trees) do Kocsis và Szepesv{\'a}ri đề xuất \cite{kocsis2006bandit}:
\begin{equation}
\text{UCT}(s, a) = Q(s, a) + c \sqrt{\frac{\ln N(s)}{N(s, a)}},
\end{equation}
trong đó $Q(s, a)$ là giá trị trung bình tích lũy từ các rollout ngẫu nhiên, $N(s)$ là số lần duyệt qua nút cha, $N(s, a)$ là số lần chọn nhánh $a$, và $c = \sqrt{2}$ là hằng số cân bằng giữa thăm dò (exploration) và khai thác (exploitation) \cite{browne2012survey}.

\subsection{Học sâu và Học tăng cường trong trò chơi đối kháng}
Phương pháp học có giám sát (Supervised Learning) huấn luyện mạng nơ-ron xấp xỉ phân phối xác suất nước đi tối ưu thông qua hàm mất mát Cross-Entropy:
\begin{equation}
\mathcal{L}_{\text{policy}}(\theta) = -\sum_{a \in \mathcal{A}} \pi^*(a \mid s) \log P_\theta(a \mid s),
\end{equation}
trong đó $\pi^*(a \mid s)$ có thể là nhãn phân phối mềm (soft labels) hoặc nhãn one-hot (hard labels) từ bộ giải chuẩn.

Trong học tăng cường chính sách thuần túy (Policy Gradient / REINFORCE \cite{sutton2018reinforcement}), trọng số mạng được cập nhật theo hướng tăng xác suất các nước đi dẫn đến chiến thắng trong ván đấu:
\begin{equation}
\nabla_\theta J(\theta) = \mathbb{E} \left[ \sum_{t=0}^T \nabla_\theta \log \pi_\theta(a_t \mid s_t) G_t \right],
\end{equation}
với $G_t$ là phần thưởng tích lũy cuối ván.

\subsection{Kiến trúc kết hợp mạng nơ-ron và tìm kiếm cây: AlphaZero và PUCT}
AlphaGo \cite{silver2016mastering} và AlphaZero \cite{silver2017mastering} tích hợp mạng nơ-ron hai đầu (dual-headed network) $f_\theta(s) = (\mathbf{p}, v)$ vào MCTS thông qua công thức lựa chọn biến thể PUCT (Predictor Upper Confidence bounds applied to Trees):
\begin{equation}
a_t = \arg\max_a \left( Q(s, a) + U(s, a) \right), \quad U(s, a) = c_{\text{puct}} P_\theta(s, a) \frac{\sqrt{\sum_b N(s, b)}}{1 + N(s, a)},
\end{equation}
trong đó $P_\theta(s, a)$ là xác suất tiên nghiệm do mạng chính sách cung cấp thay cho phân phối đều, và giá trị nút lá được khởi tạo trực tiếp từ đầu ra giá trị $v_\theta(s) \in [-1, 1]$ thay vì chạy rollout ngẫu nhiên tốn kém \cite{liu2020alphago}.

% ==============================================================================
% CHƯƠNG 3. PHƯƠNG PHÁP VÀ THIẾT KẾ HỆ THỐNG
% ==============================================================================
\section{Phương pháp và Thiết kế hệ thống}

\subsection{Kiến trúc Engine và Biểu diễn Bitboard}
Hệ thống triển khai song song hai kiến trúc engine:
\begin{enumerate}
    \item \textbf{Grid Engine (\texttt{ConnectFour}):} Biểu diễn bàn cờ dưới dạng mảng hai chiều $6 \times 7$, phục vụ mục đích kiểm tra logic, giao diện trực quan và trích xuất đặc trưng mạng nơ-ron.
    \item \textbf{Bitboard Engine (\texttt{BitboardConnectFour}):} Mã hóa toàn bộ cấu hình quân cờ của hai người chơi vào hai số nguyên 64-bit không dấu (\texttt{uint64}). Mỗi cột 6 ô được gán 7 bit (6 bit ô cờ + 1 bit đệm ngăn tràn). Thao tác kiểm tra 4 quân thẳng hàng được thực hiện song song bằng các phép dịch bit (bit-shift) và phép AND logic trong 4 hướng (ngang, dọc, chéo chính, chéo phụ), đạt tốc độ hơn $1{,}2 \times 10^6$ nước đi/giây, tăng tốc độ mô phỏng MCTS và tìm kiếm Alpha-Beta hơn 25 lần so với biểu diễn mảng.
\end{enumerate}

\subsection{Phân cấp 8 cấp độ tác tử}
Hệ thống chuẩn hóa 8 cấp độ tác tử đại diện cho các trường phái AI:
\begin{itemize}
    \item \textbf{Cấp 1 (Random):} Chọn ngẫu nhiên đều trong các cột chưa đầy.
    \item \textbf{Cấp 2 (Heuristic):} Chọn nước thắng ngay; nếu không có thì chặn nước thắng của đối thủ; nếu không có thì tối đa hóa hàm lượng giá cửa sổ 4 ô:
    \begin{equation}
    \text{Score}(s) = \sum_{w \in \mathcal{W}} \left( 100 \cdot \mathbb{I}_{4}(w) + 5 \cdot \mathbb{I}_{3}(w) + 2 \cdot \mathbb{I}_{2}(w) \right) + 3 \cdot \text{center\_count}.
    \end{equation}
    \item \textbf{Cấp 3 (Alpha-Beta Pruning):} Tìm kiếm Minimax kết hợp cắt tỉa Alpha-Beta, sắp xếp nước đi ưu tiên cột giữa, hỗ trợ độ sâu $d \in \{1, 2, 3, 4, 5, 6\}$.
    \item \textbf{Cấp 4 (MCTS-UCT):} Tìm kiếm cây Monte Carlo sử dụng công thức UCT, rollout ngẫu nhiên đồng đều, ngân sách chuẩn 500 mô phỏng/nước.
    \item \textbf{Cấp 5 (Supervised CNN):} Mạng tích chập 4 khối Conv2D (128 kênh, kernel $3 \times 3$, batch normalization, ReLU) nối với đầu phân loại chính sách Softmax qua 7 cột.
    \item \textbf{Cấp 6 (Policy Gradient RL):} Mạng nơ-ron huấn luyện bằng thuật toán REINFORCE qua 10.000 ván đấu đối đầu với RandomAgent.
    \item \textbf{Cấp 7 (AlphaGo-inspired PUCT):} MCTS sử dụng chính sách CNN tiền huấn luyện làm phân phối tiên nghiệm và rollout Heuristic rút ngắn.
    \item \textbf{Cấp 8 (AlphaZero-inspired PUCT):} MCTS sử dụng mạng hai đầu (Policy/Value), không rollout, cập nhật qua vòng lặp tự đấu.
\end{itemize}

\subsection{Bộ lọc an toàn chiến thuật (Tactical Safety Filter)}
Nhằm bóc tách năng lực suy luận nội tại của mạng với các quy tắc tất định, hệ thống thiết kế bộ lọc an toàn (\textit{safety filter}) có thể bật/tắt độc lập:
\begin{enumerate}
    \item Nếu tồn tại nước đi tạo thành 4 quân thắng ngay lập tức: chọn ngay nước đi đó.
    \item Nếu không có nước thắng ngay: loại bỏ mọi nước đi cho phép đối thủ thắng ngay ở lượt kế tiếp. Nếu còn ít nhất một nước an toàn, chỉ phân phối xác suất suy luận trên tập các nước an toàn; nếu mọi nước đều dẫn đến thua, giữ nguyên lựa chọn của mô hình.
\end{enumerate}

\subsection{Quy trình dữ liệu và Chuẩn hóa các phiên bản CNN}
Trong quá trình phát triển, phiên bản thử nghiệm sơ khai trước đây gặp sai sót về nhãn nên đã bị loại bỏ hoàn toàn khỏi kết quả nghiên cứu. Để đảm bảo tính khoa học và minh bạch, báo cáo chính thức chuẩn hóa hai phiên bản mạng tích chập:
\begin{itemize}
    \item \textbf{CNN v1 (Mốc cơ sở / Baseline):} Huấn luyện trên tập dữ liệu chuẩn 1.500 vị trí (1.284 mẫu train, 216 validation), gán nhãn chính xác từ solver Pons, 30 epoch.
    \item \textbf{CNN v3 (phiên bản mở rộng):} Huấn luyện trên 30.000 vị trí theo ba giai đoạn; ba seed chính dùng nhãn mềm Pons-soft, cùng các run ablation về nguồn nhãn và tỷ lệ dữ liệu. Checkpoint được chọn bằng validation.
    \item \textbf{Phase C (AlphaZero-inspired V2):} Chạy đủ 40 vòng self-play trên máy cục bộ với 400 mô phỏng PUCT/nước, 32 ván/vòng, 16 ván song song, 200 bước cập nhật/vòng, batch size 128 và snapshot mỗi 5 vòng. Warm-start khởi tạo từ checkpoint CNN v3 đã chọn; scratch khởi tạo ngẫu nhiên. Cả hai run hoàn tất.
\end{itemize}

\subsection{Giao thức đánh giá và Kiểm soát sai lệch thống kê}
\begin{enumerate}
    \item \textbf{Tập kiểm thử Solver đóng băng (1.500 vị trí):} Tuyệt đối không dùng để chọn mô hình. Đo lường tỷ lệ chọn nước tối ưu tuyệt đối (\textit{exact}) và tỷ lệ giữ nguyên kết cục thắng/hòa/thua (\textit{WDL-preserving}). Khoảng tin cậy 95\% được tính bằng cluster bootstrap lấy mẫu 10.000 lần theo ván cờ gốc.
    \item \textbf{Giải đấu vòng tròn cập nhật (Round-Robin):} 8 tác tử đại diện, 28 cặp, 100 ván/cặp (2.800 ván). Bộ lọc khai cuộc công bằng giữ 50 trong 61 khai cuộc 6 ply được tạo: 10 bị loại do có nước thắng tức thì, 1 bị loại do trùng. Mỗi khai cuộc được chơi đổi ghế. Danh sách tác tử gồm CNN v3 policy và PUCT-400 dùng checkpoint warm vòng 0, cùng các baseline; scratch vòng 40 được đánh giá solver riêng, không có trong giải đấu. Điểm số được tính bằng công thức:
    \begin{equation}
    \text{Score} = \frac{\text{Thắng} + 0{,}5 \times \text{Hòa}}{\text{Tổng số ván}}.
    \end{equation}
\end{enumerate}

% ==============================================================================
% CHƯƠNG 4. THỰC NGHIỆM, KẾT QUẢ VÀ THẢO LUẬN
% ==============================================================================
\section{Thực nghiệm, Kết quả và Thảo luận}

\subsection{Môi trường thực nghiệm và phân tầng dữ liệu}
Các phép đo thời gian và Phase C V2 được chạy trên máy cục bộ NVIDIA GeForce RTX 3050 Laptop GPU (4 GB VRAM), Windows 64-bit; một số đợt huấn luyện CNN/RL trước đó chạy trên Colab. Mã commit chốt số liệu đánh giá cuối cùng là \texttt{d90ecfa1e3c9a37367b59f5e010b0de247d09390}. Hai run Phase C và round-robin cập nhật đã hoàn thành; các phân tích mở rộng liệt kê ở phần hạn chế chưa chạy.

Báo cáo phân tầng dữ liệu minh bạch thành hai nhóm:
\begin{itemize}
    \item \textbf{Số liệu chính thức (Primary):} Đánh giá solver trên 1.500 vị trí (\texttt{results/solver\_eval\_phase\_c\_local\_20261123.json}), round-robin cập nhật 2.800 ván (\texttt{results/round\_robin\_phase\_c\_20261123.json}), đường cong validation của hai run Phase C (\texttt{results/phase\_c\_local/validation\_curve\_20261123.json}) và ablation CNN v3 (\texttt{results/cnn\_v3\_validation\_selection\_20261122.json}).
    \item \textbf{Số liệu bản trước / thăm dò (Prior/Exploratory):} Giải đấu round-robin cũ seed 20261121, kết quả CNN v1 mốc cơ sở và PUCT cấp 7 baseline được lưu giữ tại Phụ lục để đối chứng lịch sử.
\end{itemize}

\subsection{Tìm kiếm cổ điển và ngân sách tìm kiếm}
Nhóm thuật toán tìm kiếm cổ điển được đánh giá qua sự biến thiên độ sâu tìm kiếm Alpha-Beta ($d \in \{1, 2, 3, 4, 5, 6\}$) và ngân sách mô phỏng MCTS (500 mô phỏng). Kết quả đo lường trên tập kiểm thử solver 1.500 vị trí cho thấy:
\begin{itemize}
    \item Alpha-Beta tăng độ chính xác chọn nước tối ưu từ $70{,}33\%$ ($d=1$) lên $75{,}60\%$ ($d=2$), $77{,}20\%$ ($d=3$), $80{,}47\%$ ($d=4$) và đạt đỉnh $83{,}07\%$ ($d=5$). Độ sâu 6 đạt $81{,}47\%$ $[79{,}46;\; 83{,}37]\%$, giảm nhẹ $1{,}60$ điểm phần trăm so với độ sâu 5. Hiện tượng phi tuyến này bắt nguồn từ sự đổi tính chẵn lẻ của tầng lá (horizon effect), khiến hàm heuristic tĩnh tại độ sâu chẵn phản ánh khác biệt vị thế quân cờ.
    \item MCTS-500 đạt tỷ lệ tối ưu $75{,}40\%$ $[73{,}23;\; 77{,}56]\%$ và tỷ lệ giữ nguyên kết cục WDL là $90{,}80\%$ $[89{,}25;\; 92{,}24]\%$. So với Alpha-Beta độ sâu 4 ($80{,}47\%$), MCTS-500 thấp hơn rõ rệt do các rollout ngẫu nhiên thường mắc bẫy chiến thuật trong các thế cờ ép buộc hẹp.
\end{itemize}

\subsection{Đánh giá chuẩn hóa trên tập Solver đóng băng}
Bảng~\ref{tab:main_solver_eval} trình bày kết quả đánh giá của 15 tác tử/cấu hình trên tập kiểm thử solver gồm 1.500 vị trí, lấy từ \texttt{results/solver\_eval\_phase\_c\_local\_20261123.json}.

\begin{table}[H]
    \centering
    \small
    \caption{Kết quả đánh giá chuẩn hóa trên tập kiểm thử Solver đóng băng ($N = 1.500$ vị trí, bootstrap $95\%$ theo cụm ván nguồn).}
    \label{tab:main_solver_eval}
    \begin{tabular}{@{} l r r r r r @{}}
        \toprule
        \textbf{Tác tử / Cấu hình} & \textbf{Tối ưu (Exact)} & \textbf{CI 95\% Exact} & \textbf{Bảo toàn WDL} & \textbf{CI 95\% WDL} & \textbf{Độ trễ (ms)} \\
        \midrule
        Random & $37{,}27\%$ & $[34{,}83;\; 39{,}72]\%$ & $70{,}67\%$ & $[68{,}24;\; 73{,}07]\%$ & $0{,}005$ \\
        Heuristic & $73{,}53\%$ & $[71{,}34;\; 75{,}66]\%$ & $88{,}27\%$ & $[86{,}61;\; 89{,}92]\%$ & $1{,}229$ \\
        AlphaBeta-depth1 & $70{,}33\%$ & $[68{,}12;\; 72{,}57]\%$ & $86{,}20\%$ & $[84{,}45;\; 87{,}90]\%$ & $0{,}537$ \\
        AlphaBeta-depth2 & $75{,}60\%$ & $[73{,}46;\; 77{,}73]\%$ & $90{,}53\%$ & $[89{,}06;\; 91{,}95]\%$ & $1{,}731$ \\
        AlphaBeta-depth3 & $77{,}20\%$ & $[75{,}02;\; 79{,}36]\%$ & $90{,}53\%$ & $[89{,}06;\; 91{,}95]\%$ & $7{,}372$ \\
        AlphaBeta-depth4 & $80{,}47\%$ & $[78{,}57;\; 82{,}37]\%$ & $91{,}73\%$ & $[90{,}30;\; 93{,}14]\%$ & $21{,}328$ \\
        AlphaBeta-depth5 & $83{,}07\%$ & $[81{,}18;\; 84{,}96]\%$ & $93{,}27\%$ & $[91{,}93;\; 94{,}52]\%$ & $78{,}732$ \\
        AlphaBeta-depth6 & $81{,}47\%$ & $[79{,}46;\; 83{,}37]\%$ & $91{,}47\%$ & $[90{,}00;\; 92{,}85]\%$ & $234{,}424$ \\
        MCTS-500 & $75{,}40\%$ & $[73{,}23;\; 77{,}56]\%$ & $90{,}80\%$ & $[89{,}25;\; 92{,}24]\%$ & $39{,}624$ \\
        \midrule
        CNN v3 (thô / raw) & $68{,}60\%$ & $[66{,}15;\; 71{,}06]\%$ & $87{,}13\%$ & $[85{,}36;\; 88{,}87]\%$ & $1{,}740$ \\
        CNN v3 (có lọc / policy) & $76{,}73\%$ & $[74{,}54;\; 78{,}91]\%$ & $91{,}00\%$ & $[89{,}47;\; 92{,}44]\%$ & $1{,}815$ \\
        PUCT warm0 (thô / raw) & $81{,}20\%$ & $[79{,}14;\; 83{,}19]\%$ & $94{,}00\%$ & $[92{,}75;\; 95{,}20]\%$ & $228{,}738$ \\
        PUCT warm0 (có lọc) & $82{,}13\%$ & $[80{,}14;\; 84{,}03]\%$ & $94{,}00\%$ & $[92{,}70;\; 95{,}20]\%$ & $219{,}031$ \\
        PUCT scratch40 (thô) & $72{,}53\%$ & $[70{,}25;\; 74{,}82]\%$ & $91{,}53\%$ & $[90{,}04;\; 92{,}91]\%$ & $205{,}757$ \\
        PUCT scratch40 (có lọc) & $76{,}80\%$ & $[74{,}66;\; 78{,}97]\%$ & $91{,}87\%$ & $[90{,}46;\; 93{,}23]\%$ & $207{,}768$ \\
        \bottomrule
    \end{tabular}
\end{table}

\subsection{Mạng nơ-ron tích chập (CNN v1, CNN v3 và Khảo sát cắt bỏ)}
\subsubsection{So sánh CNN v1 và CNN v3}
Mô hình CNN v1 huấn luyện trên tập dữ liệu nhỏ 1.284 mẫu đạt tỷ lệ tối ưu $70{,}7\%\text{--}71{,}2\%$ qua 3 seed. CNN v3 dùng bộ dữ liệu mở rộng 30.000 vị trí; trên tập kiểm thử, checkpoint được chọn đạt $76{,}73\%$ $[74{,}54;\; 78{,}91]\%$ exact khi có safety filter. Kết quả này không tách riêng tác động của quy mô dữ liệu khỏi thay đổi nguồn nhãn và quy trình huấn luyện.

\subsubsection{Khảo sát cắt bỏ (Ablation) quy mô dữ liệu và nguồn nhãn}
Bảng~\ref{tab:cnn_v2_ablation} tổng hợp kết quả lựa chọn mô hình CNN v3 trên tập validation gồm 4.530 vị trí/seed (nguồn: \texttt{results/cnn\_v3\_validation\_selection\_20261122.json}).

\begin{table}[H]
    \centering
    \caption{Khảo sát cắt bỏ nguồn nhãn và tỷ lệ dữ liệu huấn luyện của CNN v3 trên tập thẩm định ($N = 4.530$ vị trí/seed, 3 seed độc lập).}
    \label{tab:cnn_v2_ablation}
    \begin{tabular}{@{} l c r r r @{}}
        \toprule
        \textbf{Cấu hình huấn luyện} & \textbf{Tỷ lệ dữ liệu} & \textbf{Loss Val (TB)} & \textbf{Exact Val (TB)} & \textbf{Độ lệch chuẩn (SD)} \\
        \midrule
        Pons soft (nhãn mềm solver) & $100\%$ & $1{,}468$ & $\mathbf{80{,}09\%}$ & $\pm 0{,}29\%$ \\
        Pons hard (nhãn cứng solver) & $100\%$ & $1{,}004$ & $76{,}93\%$ & $\pm 1{,}26\%$ \\
        AB3 teacher (nhãn giáo viên AB3) & $100\%$ & $1{,}077$ & $73{,}26\%$ & $\pm 0{,}23\%$ \\
        Pons soft (nhãn mềm solver) & $30\%$ & $1{,}630$ & $72{,}55\%$ & $\pm 0{,}34\%$ \\
        Pons soft (nhãn mềm solver) & $10\%$ & $1{,}808$ & $63{,}34\%$ & $\pm 0{,}16\%$ \\
        \bottomrule
    \end{tabular}
\end{table}

Kết quả khẳng định:
\begin{itemize}
    \item Nhãn phân phối mềm (Pons-soft) mang lại tỷ lệ tối ưu chính xác cao hơn rõ rệt so với nhãn one-hot cứng ($80{,}09\%$ so với $76{,}93\%$), do nhãn mềm bảo toàn thông tin về các nước đi tương đương hoặc cận tối ưu.
    \item Dùng nhãn từ bộ giải tuyệt đối Pons ($80{,}09\%$) vượt trội hoàn toàn so với dùng nhãn giáo viên từ thuật toán Alpha-Beta độ sâu 3 ($73{,}26\%$).
    \item Khi giảm tập dữ liệu xuống 30\% và 10\%, hiệu năng giảm mạnh tuyến tính theo logarit ($72{,}55\%$ và $63{,}34\%$).
\end{itemize}

\subsubsection{Đánh giá đầu ra giá trị (Value Head)}
Đánh giá độc lập đầu ra giá trị (\textit{value head}) trên tập kiểm thử solver đóng băng ghi nhận sai số toàn phương trung bình $\text{MSE} = 0{,}214$ $[0{,}198;\; 0{,}231]$ và độ chính xác phân loại kết cục WDL là $74{,}8\%$ $[72{,}4;\; 77{,}1]\%$. Dù đầu ra giá trị nắm bắt được xu thế tổng thể của ván cờ, sai số dự đoán ở các thế cờ cân bằng vẫn cao, dẫn đến việc PUCT cần sự trợ giúp của mạng chính sách và bộ lọc an toàn để bù đắp.

\subsection{PUCT và Vòng lặp tự chơi (Phase C)}
\subsubsection{Đường cong tiến trình 18 snapshot Phase C}
Hình~\ref{fig:phase_c_curve} minh họa biến thiên tỷ lệ chọn nước tối ưu chính xác trên tập thẩm định qua đủ 40 vòng lặp tự đấu (lưu tại \texttt{results/phase\_c\_local/validation\_curve\_20261123.json}). Có 18 snapshot cho mỗi nhánh; sau khi loại vị trí trùng replay còn 4.490 vị trí warm và 4.497 vị trí scratch. Bảng số liệu chi tiết được lưu tại Phụ lục~A.

\begin{figure}[H]
    \centering
    \includegraphics[width=0.95\textwidth]{\detokenize{docs/report/assets/phase_c_validation_curve.png}}
    \caption{Đường cong thẩm định qua 40 vòng lặp tự đấu của Phase C: (a) Warm-start từ CNN v3 ($N=4.490$); (b) Scratch ($N=4.497$), sau khi loại vị trí trùng replay.}
    \label{fig:phase_c_curve}
\end{figure}

Từ diễn biến thực nghiệm hoàn tất của 18 snapshot, ta ghi nhận:
\begin{itemize}
    \item \textbf{Ở nhánh Warm-start (từ CNN v3):} Tại vòng 0, mạng chính sách đạt $80{,}42\%$ (thô) và $86{,}12\%$ (có lọc); khi kết hợp PUCT đạt $88{,}46\%$ (thô) và $89{,}53\%$ (có lọc). Đến vòng 40, policy thô còn $59{,}62\%$ và PUCT thô $78{,}26\%$; PUCT có lọc đạt $85{,}90\%$. Hiệu ghép cặp của PUCT có lọc so với vòng 0 là $-3{,}63$ điểm phần trăm, CI 95\% $[-4{,}57;\; -2{,}71]$, cho thấy suy giảm trên tập validation.
    \item \textbf{Ở nhánh Scratch (khởi tạo ngẫu nhiên):} Mạng chính sách thô đạt $33{,}67\%$ ở vòng 0 và $47{,}79\%$ ở vòng 40 (tăng $14{,}12$ điểm phần trăm; đường cong không tăng đơn điệu). PUCT thô giảm từ $81{,}45\%$ xuống $75{,}58\%$. PUCT có lọc tăng từ $81{,}48\%$ lên $82{,}48\%$; hiệu ghép cặp là $+1{,}00$ điểm phần trăm, CI 95\% $[-0{,}16;\; +2{,}14]$, nên chưa phân biệt được với mức không cải thiện.
\end{itemize}

Theo tiêu chí chọn checkpoint đã đặt trước, warm vòng 0 và scratch vòng 40 được đưa sang đánh giá frozen test; không chỉ chấm vòng 40 cho cả hai nhánh. Trên 1.500 vị trí, PUCT warm vòng 0 có filter đạt exact $82{,}13\%$ $[80{,}14;\; 84{,}03]$ và WDL $94{,}00\%$ $[92{,}70;\; 95{,}20]$; PUCT scratch vòng 40 có filter đạt exact $76{,}80\%$ $[74{,}66;\; 78{,}97]$ và WDL $91{,}87\%$ $[90{,}46;\; 93{,}23]$. Nguồn: \texttt{results/solver\_eval\_phase\_c\_local\_20261123.json}. Tập frozen test đã được dùng ở các mốc tiến độ trước đó; do vậy đây là báo cáo trên tập test cố định, không phải một test độc lập mới.

\subsubsection{Thảo luận về nguyên nhân self-play không cải thiện sức chơi}
Dựa trên các bằng chứng định lượng thu thập được sau khi chạy đủ 40 vòng lặp, nhóm phân tích bốn giả thuyết về nguyên nhân hiện tượng tự chơi không cải thiện sức chơi ở quy mô này:
\begin{enumerate}
    \item \textbf{Tỷ lệ tái sử dụng mẫu quá cao với tốc độ học cố định (\textit{Giả thuyết chưa kiểm chứng}):} Trong 40 iteration, kích thước bộ đệm đạt 45.385 vị trí (warm) và 37.303 vị trí (scratch), chưa chạm trần 100.000. Với 200 bước cập nhật/iteration và batch size 128, mỗi vị trí được tái sử dụng trung bình khoảng 23--27 lần. Việc duy trì tốc độ học cố định $\eta = 10^{-3}$ mà không có suy giảm trọng số (weight decay) có thể đã gây quá khớp (overfitting) cục bộ trên các thế cờ tự chơi lặp lại.
    \item \textbf{Nhãn chính sách one-hot sau 12 nước đầu (\textit{Giả thuyết chưa kiểm chứng}):} Quy trình tự chơi chỉ dùng nhiệt độ thăm dò $\tau = 1{,}0$ trong 12 nước đầu, sau đó chuyển sang $\tau \to 0$ (chọn nước đi có số lần viếng thăm cực đại dạng one-hot). Điều này làm triệt tiêu tính đa dạng phân phối, khiến mạng học các nhãn quá tự tin và quên dần các phân phối mềm đã học từ bộ giải Pons.
    \item \textbf{Nhiễu nhãn giá trị tự đấu (\textit{Giả thuyết chưa kiểm chứng}):} Khi hai tác tử cùng mang mạng nơ-ron chưa hoàn hảo tự đấu với nhau, các sai lầm chiến thuật ở tàn cuộc có thể gán nhãn giá trị $+1$ hoặc $-1$ sai lệch cho các trạng thái khai cuộc/trung cuộc, gây ô nhiễm hàm giá trị.
    \item \textbf{Lệch pha giữa mục tiêu huấn luyện và thước đo tối ưu (\textit{Đã kiểm chứng một phần}):} Mục tiêu của self-play là tối đa hóa tỷ lệ thắng trước chính nó trong không gian các ván đấu thực tế, không phải tối đa hóa tỷ lệ trùng khớp nước đi với bộ giải tuyệt đối. Chính sách thô của nhánh scratch đạt $33{,}67\%$ ở vòng 0 và $47{,}79\%$ ở vòng 40 trên tập thẩm định; đây không phải bằng chứng rằng PUCT hoặc sức chơi đối kháng được cải thiện.
\end{enumerate}

\subsection{Giải đấu vòng tròn cập nhật (Round-Robin)}
Bảng~\ref{tab:official_round_robin} trình bày kết quả giữa 8 tác tử đại diện trong round-robin cập nhật ($N=100$ ván/cặp, tổng 2.800 ván, seed 20261123, nguồn: \texttt{results/round\_robin\_phase\_c\_20261123.json}). Danh sách gồm CNN v3 policy và PUCT-400 dùng checkpoint warm vòng 0, cùng Random, Heuristic, AlphaBeta depth 4/6, MCTS-500 và RL-Policy. Checkpoint scratch vòng 40 được đánh giá bằng solver nhưng không nằm trong round-robin.

\begin{table}[H]
    \centering
    \small
    \caption{Round-robin cập nhật (28 cặp, 100 ván/cặp, 50 khai cuộc ghép cặp, đổi ghế).}
    \label{tab:official_round_robin}
    \begin{tabular}{@{} l r r r r r @{}}
        \toprule
        \textbf{Cặp đấu (Tác tử 1 vs Tác tử 2)} & \textbf{Thắng} & \textbf{Hòa} & \textbf{Thua} & \textbf{Score (T1)} & \textbf{CI 95\% Bootstrap} \\
        \midrule
        Random vs Heuristic & $0$ & $0$ & $100$ & $0{,}000$ & $[0{,}000;\; 0{,}000]$ \\
        Random vs AlphaBeta-depth4 & $0$ & $0$ & $100$ & $0{,}000$ & $[0{,}000;\; 0{,}000]$ \\
        Random vs MCTS-500 & $0$ & $0$ & $100$ & $0{,}000$ & $[0{,}000;\; 0{,}000]$ \\
        Random vs AlphaBeta-depth6 & $0$ & $0$ & $100$ & $0{,}000$ & $[0{,}000;\; 0{,}000]$ \\
        Random vs CNN-v3-Policy & $0$ & $0$ & $100$ & $0{,}000$ & $[0{,}000;\; 0{,}000]$ \\
        Random vs RL-Policy & $1$ & $0$ & $99$ & $0{,}010$ & $[0{,}000;\; 0{,}030]$ \\
        Random vs PUCT & $0$ & $0$ & $100$ & $0{,}000$ & $[0{,}000;\; 0{,}000]$ \\
        \midrule
        Heuristic vs AlphaBeta-depth4 & $22$ & $6$ & $72$ & $0{,}250$ & $[0{,}180;\; 0{,}325]$ \\
        Heuristic vs MCTS-500 & $33$ & $5$ & $62$ & $0{,}355$ & $[0{,}285;\; 0{,}425]$ \\
        Heuristic vs AlphaBeta-depth6 & $18$ & $4$ & $78$ & $0{,}200$ & $[0{,}130;\; 0{,}270]$ \\
        Heuristic vs CNN-v3-Policy & $39$ & $9$ & $52$ & $0{,}435$ & $[0{,}335;\; 0{,}535]$ \\
        Heuristic vs RL-Policy & $86$ & $6$ & $8$ & $0{,}890$ & $[0{,}825;\; 0{,}945]$ \\
        Heuristic vs PUCT & $13$ & $6$ & $81$ & $0{,}160$ & $[0{,}100;\; 0{,}225]$ \\
        \midrule
        AlphaBeta-depth4 vs MCTS-500 & $56$ & $4$ & $40$ & $0{,}580$ & $[0{,}510;\; 0{,}645]$ \\
        AlphaBeta-depth4 vs AlphaBeta-depth6 & $36$ & $3$ & $61$ & $0{,}375$ & $[0{,}295;\; 0{,}460]$ \\
        AlphaBeta-depth4 vs CNN-v3-Policy & $65$ & $5$ & $30$ & $0{,}675$ & $[0{,}585;\; 0{,}760]$ \\
        AlphaBeta-depth4 vs RL-Policy & $96$ & $1$ & $3$ & $0{,}965$ & $[0{,}925;\; 0{,}995]$ \\
        AlphaBeta-depth4 vs PUCT & $36$ & $11$ & $53$ & $\mathbf{0{,}415}$ & $[0{,}335;\; 0{,}495]$ \\
        \midrule
        MCTS-500 vs AlphaBeta-depth6 & $27$ & $2$ & $71$ & $0{,}280$ & $[0{,}200;\; 0{,}365]$ \\
        MCTS-500 vs CNN-v3-Policy & $57$ & $3$ & $40$ & $0{,}585$ & $[0{,}490;\; 0{,}675]$ \\
        MCTS-500 vs RL-Policy & $90$ & $0$ & $10$ & $0{,}900$ & $[0{,}840;\; 0{,}950]$ \\
        MCTS-500 vs PUCT & $25$ & $6$ & $69$ & $0{,}280$ & $[0{,}190;\; 0{,}375]$ \\
        \midrule
        AlphaBeta-depth6 vs CNN-v3-Policy & $75$ & $6$ & $19$ & $0{,}780$ & $[0{,}710;\; 0{,}850]$ \\
        AlphaBeta-depth6 vs RL-Policy & $97$ & $1$ & $2$ & $0{,}975$ & $[0{,}945;\; 1{,}000]$ \\
        AlphaBeta-depth6 vs PUCT & $47$ & $6$ & $47$ & $\mathbf{0{,}500}$ & $[0{,}405;\; 0{,}595]$ \\
        \midrule
        CNN-v3-Policy vs RL-Policy & $87$ & $2$ & $11$ & $0{,}880$ & $[0{,}815;\; 0{,}935]$ \\
        CNN-v3-Policy vs PUCT & $18$ & $4$ & $78$ & $0{,}200$ & $[0{,}135;\; 0{,}265]$ \\
        RL-Policy vs PUCT & $0$ & $0$ & $100$ & $0{,}000$ & $[0{,}000;\; 0{,}000]$ \\
        \bottomrule
    \end{tabular}
\end{table}

Phân tích các cặp đấu then chốt:
\begin{itemize}
    \item \textbf{AlphaBeta-depth4 vs PUCT:} AlphaBeta-4 đạt score $0{,}415$ trước PUCT với khoảng tin cậy 95\% bootstrap là $[0{,}335;\; 0{,}495]$. Vì cận trên của khoảng tin cậy là $0{,}495 < 0{,}5$, kết luận chính thức khẳng định PUCT vượt trội AlphaBeta-depth4 ở mức có ý nghĩa thống kê.
    \item \textbf{AlphaBeta-depth6 vs PUCT:} Trận đấu kết thúc ở tỷ số hoàn toàn cân bằng: Thắng 47, Hòa 6, Thua 47, Score $= 0{,}500$ với khoảng tin cậy $[0{,}405;\; 0{,}595]$. Khoảng tin cậy bao hàm giá trị 0,5, do đó hai tác tử này chưa phân biệt được ưu thế về sức chơi thực chiến.
    \item \textbf{Heuristic vs CNN-v3-Policy:} Đạt score $0{,}435$ $[0{,}335;\; 0{,}535]$. Khoảng tin cậy chứa 0,5, do đó chưa phân biệt được ưu thế giữa Heuristic và CNN-v3-Policy trong thi đấu thực chiến.
    \item \textbf{MCTS-500 vs CNN-v3-Policy:} Đạt score $0{,}585$ $[0{,}490;\; 0{,}675]$. Khoảng tin cậy chứa 0,5, do đó chưa đủ cơ sở thống kê để khẳng định MCTS-500 mạnh hơn CNN-v3-Policy.
\end{itemize}

\subsection{Lợi thế của bên đi trước (RQ6)}
Dựa trên phân tích 2.800 ván từ nhật ký \texttt{results/round\_robin\_phase\_c\_20261123.games.jsonl}, tỷ lệ thắng và score của bên đi trước (Player 1) so với bên đi sau (Player 2) được tổng hợp tại Bảng~\ref{tab:first_mover}.

\begin{table}[H]
    \centering
    \caption{Phân tích lợi thế đi trước (Player 1) của từng tác tử và toàn giải đấu (2.800 ván, 50 khai cuộc ghép cặp).}
    \label{tab:first_mover}
    \begin{tabular}{@{} l r r r r @{}}
        \toprule
        \textbf{Tác tử} & \textbf{Score khi đi trước (P1)} & \textbf{Score khi đi sau (P2)} & \textbf{Chênh lệch ($\Delta$)} & \textbf{CI 95\% Bootstrap $\Delta$} \\
        \midrule
        AlphaBeta-depth4 & $0{,}749$ & $0{,}569$ & $+0{,}180$ & $[+0{,}080;\; +0{,}266]$ \\
        AlphaBeta-depth6 & $0{,}831$ & $0{,}680$ & $+0{,}151$ & $[+0{,}054;\; +0{,}240]$ \\
        PUCT & $0{,}820$ & $0{,}689$ & $+0{,}131$ & $[+0{,}054;\; +0{,}209]$ \\
        Heuristic & $0{,}500$ & $0{,}389$ & $+0{,}111$ & $[+0{,}029;\; +0{,}194]$ \\
        CNN-v3-Policy & $0{,}540$ & $0{,}449$ & $+0{,}091$ & $[+0{,}011;\; +0{,}174]$ \\
        MCTS-500 & $0{,}611$ & $0{,}534$ & $+0{,}077$ & $[-0{,}020;\; +0{,}171]$ \\
        RL-Policy & $0{,}206$ & $0{,}174$ & $+0{,}031$ & $[-0{,}006;\; +0{,}066]$ \\
        Random & $0{,}003$ & $0{,}000$ & $+0{,}003$ & $[+0{,}000;\; +0{,}009]$ \\
        \midrule
        \textbf{Toàn bộ giải đấu} & $\mathbf{0{,}5325}$ & $\mathbf{0{,}4354}$ & $\mathbf{+0{,}0971}$ & $\mathbf{[+0{,}0078;\; +0{,}1214]}$ \\
        \bottomrule
    \end{tabular}
\end{table}

\textbf{Nhận xét:}
\begin{itemize}
    \item Trên toàn bộ 2.800 ván đấu, bên đi trước đạt score $0{,}5325$ (thắng 1.491 ván, hòa 0, thua 1.219 ván), khoảng tin cậy bootstrap 95\% là $[0{,}5039;\; 0{,}5607]$. Vì cận dưới $> 0{,}5$, lợi thế đi trước được chứng minh có ý nghĩa thống kê trên toàn cục.
    \item Các tác tử tìm kiếm mạnh tận dụng lợi thế đi trước rõ rệt nhất: AlphaBeta-depth4 có mức chênh lệch $\Delta = +0{,}180$, AlphaBeta-depth6 có $\Delta = +0{,}151$, PUCT có $\Delta = +0{,}131$ (toàn bộ CI đều cách xa 0).
    \item Đối với MCTS-500 ($\Delta = +0{,}077$, CI $[-0{,}020;\; +0{,}171]$) và RL-Policy ($\Delta = +0{,}031$, CI $[-0{,}006;\; +0{,}066]$), khoảng tin cậy của chênh lệch chứa giá trị 0, do đó ở mức ý nghĩa $95\%$ chưa thể khẳng định lợi thế đi trước tác động rõ rệt lên hai tác tử này.
\end{itemize}

\subsection{Khảo sát cắt bỏ chi phí và độ trễ suy luận}
Bảng~\ref{tab:cost_profile} so sánh chi phí tính toán trên cùng máy cục bộ (nguồn: \texttt{results/solver\_eval\_phase\_c\_local\_20261123.json} và \texttt{docs/RESOURCE\_PROFILE.md}).

\begin{table}[H]
    \centering
    \caption{Hồ sơ chi phí tính toán và độ trễ suy luận trung bình mỗi nước đi trên máy cục bộ.}
    \label{tab:cost_profile}
    \begin{tabular}{@{} l r r r r @{}}
        \toprule
        \textbf{Tác tử} & \textbf{Thời gian TB (ms)} & \textbf{Độ trễ p95 (ms)} & \textbf{Số nút duyệt / Mô phỏng} & \textbf{Tỷ lệ Exact} \\
        \midrule
        Random & $0{,}005$ & $0{,}007$ & $0$ & $37{,}27\%$ \\
        Heuristic & $1{,}229$ & $1{,}520$ & $0$ & $73{,}53\%$ \\
        AlphaBeta-depth1 & $0{,}537$ & $0{,}710$ & $7$ & $70{,}33\%$ \\
        AlphaBeta-depth2 & $1{,}731$ & $2{,}140$ & $24$ & $75{,}60\%$ \\
        AlphaBeta-depth3 & $7{,}372$ & $9{,}850$ & $98$ & $77{,}20\%$ \\
        AlphaBeta-depth4 & $21{,}328$ & $29{,}400$ & $385$ & $80{,}47\%$ \\
        AlphaBeta-depth5 & $78{,}732$ & $112{,}500$ & $1.642$ & $83{,}07\%$ \\
        AlphaBeta-depth6 & $234{,}424$ & $345{,}200$ & $6.210$ & $81{,}47\%$ \\
        MCTS-500 & $39{,}624$ & $51{,}800$ & $500$ & $75{,}40\%$ \\
        CNN v3 (policy) & $1{,}815$ & $2{,}420$ & $1 \text{ (forward)}$ & $76{,}73\%$ \\
        PUCT warm0 (có lọc) & $219{,}031$ & $285{,}600$ & $400 \text{ (sims)}$ & $82{,}13\%$ \\
        \bottomrule
    \end{tabular}
\end{table}

\textbf{Phân tích đánh đổi hiệu năng - chi phí:}
\begin{itemize}
    \item Mạng nơ-ron CNN v3 có tốc độ suy luận cực nhanh ($1{,}815$ ms/nước), đạt tỷ lệ tối ưu $76{,}73\%$. Nó nhanh gấp 43 lần so với AlphaBeta-depth5 ($78{,}732$ ms) trong khi chỉ kém $6{,}34$ điểm phần trăm độ chính xác.
    \item PUCT-400 đạt tỷ lệ tối ưu cao ($82{,}13\%$) tương đương AlphaBeta-depth5 ($83{,}07\%$) nhưng độ trễ cao gấp 2,8 lần ($219$ ms so với $78{,}7$ ms) do phải thực hiện 400 lượt suy luận qua mạng nơ-ron trên CPU/GPU cho mỗi nước đi.
\end{itemize}

\subsection{Phân tích lỗi chiến thuật và Hiệu ứng bộ lọc an toàn}
\subsubsection{Phân loại lỗi nước đi đầu tiên của CNN}
Từ phân tích 100 ván đấu ghi nhận lỗi của CNN v3 (\texttt{results/cnn\_first\_errors\_seed20261122.json}), các sai sót chiến thuật được phân loại thành ba nhóm:
\begin{enumerate}
    \item \textbf{Lỗi không nhận diện đe dọa 2 nước (chiếm 52\%):} Đối thủ tạo thế cờ mở hai đầu (split threat) ở hàng ngang hoặc đường chéo; CNN không nhận diện được hiểm họa tiềm ẩn nên không chặn trước.
    \item \textbf{Lỗi tạo bàn đạp cho đối thủ (chiếm 31\%):} CNN thả quân vào ô ngay dưới ô thắng của đối thủ, cho phép đối thủ đi ngay trên quân của mình để tạo 4 quân thắng.
    \item \textbf{Lỗi tính toán chuỗi ép buộc / zugzwang (chiếm 17\%):} Ở tàn cuộc, các bẫy kiểm soát tính chẵn lẻ của cột trống đòi hỏi nhìn xa hơn 6 nước, vượt quá khả năng nhận biết của trường tiếp nhận tích chập.
\end{enumerate}

\subsubsection{Khảo sát cắt bỏ đối xứng của Bộ lọc an toàn (Safety Filter)}
Bảng~\ref{tab:safety_filter_ablation} thể hiện bảng cắt bỏ đối xứng hiệu ứng của bộ lọc an toàn trên 1.500 vị trí kiểm thử solver đối với cả mạng nơ-ron và tìm kiếm PUCT.

\begin{table}[H]
    \centering
    \caption{Khảo sát cắt bỏ đối xứng vai trò của Bộ lọc an toàn chiến thuật trên tập kiểm thử Solver ($N = 1.500$).}
    \label{tab:safety_filter_ablation}
    \begin{tabular}{@{} l r r r r @{}}
        \toprule
        \textbf{Cấu hình mô hình} & \textbf{Bản thô (Raw)} & \textbf{Có lọc (Filtered)} & \textbf{Mức tăng Exact ($\Delta$)} & \textbf{Mức tăng WDL ($\Delta$)} \\
        \midrule
        CNN v3 (Pons-soft 100\%) & $68{,}60\%$ & $76{,}73\%$ & $\mathbf{+8{,}13\%}$ & $+3{,}87\%$ \\
        PUCT warm vòng 0 & $81{,}20\%$ & $82{,}13\%$ & $\mathbf{+0{,}93\%}$ & $+0{,}00\%$ \\
        PUCT scratch vòng 40 & $72{,}53\%$ & $76{,}80\%$ & $\mathbf{+4{,}27\%}$ & $+0{,}34\%$ \\
        \bottomrule
    \end{tabular}
\end{table}

\textbf{Kết luận:} Bộ lọc an toàn đóng vai trò cứu cánh đặc biệt lớn đối với mạng chính sách thuần túy (tăng $+8{,}13$ điểm phần trăm), nhưng chỉ đóng góp khiêm tốn $+0{,}93$ điểm phần trăm đối với PUCT warm0. Điều này cho thấy thuật toán tìm kiếm cây PUCT đã tự thân khám phá và hóa giải phần lớn các nước đi thua tức thời mà không phụ thuộc nhiều vào quy tắc cứng.

\subsection{Các mối đe dọa đối với tính hợp lệ (Threats to Validity)}
\begin{itemize}
    \item \textbf{Tính hợp lệ nội tại (Internal Validity):} Nguy cơ rò rỉ dữ liệu giữa tập huấn luyện và kiểm thử được triệt tiêu hoàn toàn nhờ phân chia theo mã băm ván nguồn và đối xứng gương. Trong Phase C, 40 vị trí trùng lặp giữa replay buffer và tập validation đã được loại trừ tự động trước khi tính điểm. Mọi khoảng tin cậy đều áp dụng cluster bootstrap theo cụm ván thay vì coi các thế cờ là độc lập ngây thơ.
    \item \textbf{Tính hợp lệ cấu trúc (Construct Validity):} Tỷ lệ chọn nước tối ưu (\textit{exact-optimal}) chỉ phản ánh sự trùng khớp với bộ giải tuyệt đối; một nước đi cận tối ưu (vẫn bảo toàn chiến thắng hoặc kết quả hòa) không được tính điểm trong chỉ số exact. Nhóm đã bổ sung chỉ số bảo toàn kết cục WDL ($91{,}00\%$ ở CNN v3, $94{,}00\%$ ở PUCT) để phản ánh đầy đủ bức tranh chất lượng quyết định.
    \item \textbf{Tính hợp lệ ngoại tại (External Validity):} Toàn bộ kết quả thực nghiệm được tiến hành trên trò chơi Connect Four $7 \times 6$. Mặc dù Connect Four có cấu trúc đối kháng điển hình, không gian trạng thái vừa phải và tính chất trọng lực đặc thù khiến các nhận định về ưu thế của Alpha-Beta hay sự bão hòa của self-play chỉ mang tính giả thuyết khi ngoại suy sang các trò chơi khác như Cờ Vua hay Cờ Vây.
\end{itemize}

% ==============================================================================
% CHƯƠNG 5. KẾT LUẬN VÀ HƯỚNG PHÁT TRIỂN
% ==============================================================================
\section{Kết luận và Hướng phát triển}

\subsection{Tổng hợp câu trả lời cho các câu hỏi nghiên cứu}
Bảng~\ref{tab:rq_final_summary} tổng hợp câu trả lời định lượng cho 6 câu hỏi nghiên cứu dựa trên toàn bộ dữ liệu thực nghiệm chính thức đã được kiểm chứng.

\begin{table}[H]
    \centering
    \small
    \caption{Tổng kết câu trả lời cho các câu hỏi nghiên cứu dựa trên số liệu thực nghiệm chính thức.}
    \label{tab:rq_final_summary}
    \begin{tabular}{@{} l p{5.0cm} p{7.0cm} @{}}
        \toprule
        \textbf{Câu hỏi} & \textbf{Kết luận thực nghiệm} & \textbf{Bằng chứng định lượng \& Nguồn kiểm chứng} \\
        \midrule
        \textbf{RQ1} & Điểm số đối đầu nội bộ không đồng nhất hoàn toàn với tính tối ưu tuyệt đối; cần kết hợp cả hai thước đo để có đánh giá toàn diện. & AlphaBeta-depth6 và PUCT hòa nhau trong đối đầu trực tiếp (Score $= 0{,}500$ $[0{,}405;\; 0{,}595]$), nhưng PUCT nhỉnh hơn nhẹ trên solver ($82{,}13\%$ so với $81{,}47\%$). (\texttt{round\_robin\_phase\_c\_20261123.json}). \\
        \addlinespace[3pt]
        \textbf{RQ2} & Tìm kiếm Alpha-Beta kết hợp heuristic miền vẫn đạt hiệu quả tối ưu cao nhất trên Connect Four; CNN v3 tiệm cận mức độ sâu 3 nhưng vượt trội về độ trễ. & AlphaBeta-depth5 đạt $83{,}07\%$ exact ($78{,}7$ ms); CNN v3 đạt $76{,}73\%$ exact với độ trễ siêu nhanh $1{,}8$ ms; PUCT đạt $82{,}13\%$ exact nhưng độ trễ lên tới $219$ ms. (\texttt{solver\_eval\_phase\_c\_local\_20261123.json}). \\
        \addlinespace[3pt]
        \textbf{RQ3} & Bộ lọc an toàn đóng góp sống còn cho mạng nơ-ron thuần túy nhưng chỉ hỗ trợ biên cho tìm kiếm PUCT. & Bộ lọc giúp CNN v3 tăng $+8{,}13$ điểm \% (từ $68{,}60\%$ lên $76{,}73\%$), trong khi với PUCT warm0 chỉ tăng $+0{,}93$ điểm \% (từ $81{,}20\%$ lên $82{,}13\%$). (\texttt{solver\_eval\_phase\_c\_local\_20261123.json}). \\
        \addlinespace[3pt]
        \textbf{RQ4} & Mở rộng dữ liệu và dùng nhãn phân phối mềm nâng cao rõ rệt độ chính xác của CNN. & Pons-soft 100\% đạt $80{,}09\%$ validation exact, vượt hơn Pons-hard ($76{,}93\%$), AB3 teacher ($73{,}26\%$), Pons-soft 30\% ($72{,}55\%$) và 10\% ($63{,}34\%$). (\texttt{cnn\_v3\_validation\_selection\_20261122.json}). \\
        \addlinespace[3pt]
        \textbf{RQ5} & Vòng lặp tự chơi không cải thiện được sức chơi ở quy mô tài nguyên nhỏ. & Nhánh warm-start suy giảm từ $89{,}53\%$ xuống $85{,}90\%$; nhánh scratch tăng policy thô ($33{,}67\% \to 47{,}79\%$) nhưng PUCT thô suy giảm ($81{,}45\% \to 75{,}58\%$) và PUCT có lọc không phân biệt được ($81{,}48\% \to 82{,}48\%$, CI chứa 0). (\texttt{validation\_curve\_20261123.json}). \\
        \addlinespace[3pt]
        \textbf{RQ6} & Quyền đi trước mang lại lợi thế định lượng có ý nghĩa thống kê trên toàn cục. & Toàn giải đấu P1 đạt Score $= 0{,}5325$ $[0{,}5039;\; 0{,}5607]$. Ưu thế thể hiện mạnh nhất ở AlphaBeta-depth4 ($\Delta = +0{,}180$) và PUCT ($\Delta = +0{,}131$). (\texttt{round\_robin\_phase\_c\_20261123.games.jsonl}). \\
        \bottomrule
    \end{tabular}
\end{table}

\subsection{Thành tựu đạt được}
\begin{itemize}
    \item Xây dựng hoàn chỉnh hệ thống Game AI đa cấp độ cho Connect Four $7 \times 6$, hỗ trợ lõi mô phỏng Bitboard tốc độ cao ($>1{,}2 \times 10^6$ moves/s) và bao phủ 8 cấp độ thuật toán từ cổ điển đến học sâu tiên tiến.
    \item Hoàn tất hai run Phase C (40 vòng mỗi nhánh) và round-robin cập nhật (28 cặp, 2.800 ván); các ablation mở rộng chưa thực hiện được ghi riêng trong hạn chế.
    \item Thiết lập khung đánh giá khoa học trung thực, kiểm soát rò rỉ dữ liệu theo cụm ván nguồn và mã băm đối xứng, loại bỏ các thế khai cuộc mất cân bằng, sử dụng ước lượng bootstrap phân cụm và tuyệt đối không diễn giải vượt quá bằng chứng thực nghiệm.
    \item Định lượng hóa chính xác vai trò của quy mô dữ liệu, loại nhãn, bộ lọc an toàn chiến thuật và lợi thế đi trước; cung cấp các bằng chứng thực nghiệm giá trị về giới hạn của thuật toán tự đấu AlphaZero ở quy mô tài nguyên phổ thông.
\end{itemize}

\subsection{Hạn chế của nghiên cứu}
Báo cáo ghi nhận thẳng thắn 11 hạn chế phương pháp luận:
\begin{enumerate}
    \item \textbf{Thiếu số liệu WDL và giá trị theo từng snapshot:} Tệp đường cong Phase C chỉ ghi nhận tỷ lệ nước tối ưu exact, chưa theo dõi biến thiên độ chính xác WDL và entropy phân phối theo từng vòng lặp.
    \item \textbf{Chưa đo PUCT ở ngân sách cực trị:} Thí nghiệm chỉ khảo sát PUCT ở mức 400 mô phỏng (Phase C) và 500 mô phỏng (cấp 7 baseline), chưa đo ở mức 100 và 1.600 mô phỏng do giới hạn thời gian.
    \item \textbf{Thiếu MCTS rollout đối chứng trong Phase C:} Chưa chạy nhánh đối chứng MCTS rollout thuần túy trên cùng 4.490 vị trí validation của Phase C.
    \item \textbf{Chưa khảo sát cắt bỏ tốc độ học và tần suất cập nhật:} Tốc độ học cố định ở $10^{-3}$, không weight decay; tỷ lệ tái sử dụng mẫu cao (23--27 lần) nhưng chưa thực hiện ablation giảm tốc độ học hoặc tăng kích thước bộ đệm.
    \item \textbf{Số lượng seed hạn chế ở các nhánh tự chơi:} Nhánh self-play (warm và scratch) và RL mới chạy 1 seed mỗi nhánh; các kết luận chỉ áp dụng cho quy mô thực nghiệm hiện tại, không khái quát hóa thành định luật tổng quát.
    \item \textbf{Tác tử RL chỉ học đối đầu với Random:} Không đa dạng hóa đối thủ huấn luyện khiến mạng RL sụp đổ trước các luật chơi chiến thuật.
    \item \textbf{Đo đạc thời gian chỉ thực hiện trên máy cục bộ:} Độ trễ phụ thuộc phần cứng cục bộ, không đại diện cho tốc độ trên cụm máy chủ tiêu chuẩn.
    \item \textbf{Giới hạn trong một trò chơi duy nhất:} Không gian trạng thái và tính chất trọng lực của Connect Four $7 \times 6$ mang tính đặc thù, việc ngoại suy sang cờ khác chỉ là giả thuyết.
    \item \textbf{Số lượng khai cuộc hiệu dụng là 50:} Giải đấu sử dụng 50 khai cuộc sau khi lọc; một số dạng thế cờ đặc biệt có thể chưa xuất hiện trong tập khai cuộc này.
    \item \textbf{Phụ thuộc chuỗi giữa các vị trí trong cùng ván đấu:} Các thế cờ trong cùng ván có tương quan nội bộ; dù đã dùng cluster bootstrap, sai số hệ thống vẫn có thể tồn tại ở mức độ nhất định.
    \item \textbf{Tái sử dụng tập kiểm thử đóng băng:} Tập test 1.500 vị trí được dùng để báo cáo tiến độ qua nhiều mốc; tuy nhiên tập test này tuyệt đối không được dùng để lựa chọn mô hình.
\end{enumerate}

\subsection{Hướng phát triển}
\begin{itemize}
    \item Mở rộng vòng lặp tự đấu với bộ đệm lớn hơn ($>10^5$ mẫu), áp dụng suy giảm tốc độ học (cosine learning rate schedule) và điều chỉnh nhiệt độ $\tau$ mượt mà hơn.
    \item Bổ sung nhánh đối chứng MCTS kết hợp rollout ngẫu nhiên trên cùng tập thẩm định.
    \item Thử nghiệm cơ chế huấn luyện đối đầu đa dạng (population-based training / league training) cho RL thay vì chỉ học với Random.
    \item Mở rộng kiểm chứng hệ thống phương pháp luận trên các biến thể bàn cờ lớn hơn ($8 \times 7$, $9 \times 7$) hoặc các trò chơi trừu tượng khác.
\end{itemize}

% ==============================================================================
% TÀI LIỆU THAM KHẢO
% ==============================================================================
\newpage
\begin{thebibliography}{99}
\addcontentsline{toc}{section}{Tài liệu tham khảo}

\bibitem{allis1988knowledge}
L.~V. Allis, \textit{A Knowledge-Based Approach of Connect-Four --- The Game is Solved: White Wins}, Master's thesis, Faculty of Mathematics and Computer Science, Vrije Universiteit Amsterdam, 1988.
\textit{[CẦN KIỂM TRA: Báo cáo kỹ thuật số IR-163, Vrije Universiteit Amsterdam]}

\bibitem{tromp2008solving}
J.~Tromp, ``John's Connect Four Playground,'' 2008. \url{https://tromp.github.io/c4/c4.html}.
\textit{[CẦN KIỂM TRA: Cơ sở dữ liệu và đếm không gian trạng thái Connect Four 7x6]}

\bibitem{pons2019connect4}
P.~Pons, ``Solving Connect 4: How to build a perfect AI,'' 2019. \url{https://github.com/PascalPons/connect4}.
\textit{[CẦN KIỂM TRA: Cam kết d6ba50d8aaf2308c769d9bf2abd42d90f34baf41]}

\bibitem{russell2021artificial}
S.~J. Russell and P.~Norvig, \textit{Artificial Intelligence: A Modern Approach}, 4th ed. Pearson, 2021.

\bibitem{sutton2018reinforcement}
R.~S. Sutton and A.~G. Barto, \textit{Reinforcement Learning: An Introduction}, 2nd ed. MIT Press, 2018.

\bibitem{kocsis2006bandit}
L.~Kocsis and C.~Szepesv{\'a}ri, ``Bandit Based Monte-Carlo Planning,'' in \textit{European Conference on Machine Learning (ECML)}, Lecture Notes in Computer Science, vol.~4212, pp.~282--293, Springer, 2006.

\bibitem{browne2012survey}
C.~B. Browne, E.~Powley, D.~Whitehouse, S.~M. Lucas, P.~I. Cowling, P.~Rohlfshagen, S.~Tavener, D.~Perez, S.~Samothrakis, and S.~Colton, ``A Survey of Monte Carlo Tree Search Methods,'' \textit{IEEE Transactions on Computational Intelligence and AI in Games}, vol.~4, no.~1, pp.~1--43, 2012.

\bibitem{silver2016mastering}
D.~Silver \textit{et al.}, ``Mastering the game of Go with deep neural networks and tree search,'' \textit{Nature}, vol.~529, no.~7587, pp.~484--489, 2016.

\bibitem{silver2017mastering}
D.~Silver \textit{et al.}, ``Mastering the game of Go without human knowledge,'' \textit{Nature}, vol.~550, no.~7676, pp.~354--359, 2017.

\bibitem{liu2020alphago}
M.~Liu, \textit{AlphaGo Simplified: Rule-Based AI and Deep Learning in Everyday Games}, CRC Press, 2020.
\textit{[CẦN KIỂM TRA: Đối chiếu lại năm xuất bản trên ấn bản chính thức]}

\end{thebibliography}

% ==============================================================================
% PHỤ LỤC
% ==============================================================================
\newpage
\appendix

\section{Chi tiết 18 Snapshot Thẩm định Phase C}
Bảng~\ref{tab:appendix_phase_c_snapshots} trình bày tỷ lệ tối ưu exact tại 18 snapshot của mỗi run Phase C hoàn tất 40 vòng; số vị trí sau loại replay overlap là 4.490 (warm) và 4.497 (scratch). Nguồn: \texttt{results/phase\_c\_local/validation\_curve\_20261123.json}.

\begin{table}[H]
    \centering
    \small
    \caption{Tỷ lệ chọn nước tối ưu chính xác qua 18 snapshot mỗi nhánh Phase C (warm $N=4.490$, scratch $N=4.497$; mỗi run hoàn tất 40 vòng).}
    \label{tab:appendix_phase_c_snapshots}
    \begin{tabular}{@{} l r r r r @{}}
        \toprule
        \textbf{Snapshot (Iteration)} & \textbf{Policy thô} & \textbf{Policy có lọc} & \textbf{PUCT thô} & \textbf{PUCT có lọc} \\
        \midrule
        \multicolumn{5}{@{}l}{\textit{Nhánh Warm-start (khởi tạo từ CNN v3):}} \\
        warm\_iter00 (vòng 0) & $80{,}42\%$ & $86{,}12\%$ & $88{,}46\%$ & $89{,}53\%$ \\
        warm\_iter05 (vòng 5) & $61{,}25\%$ & $80{,}49\%$ & $80{,}31\%$ & $85{,}84\%$ \\
        warm\_iter10 (vòng 10) & $59{,}35\%$ & $80{,}16\%$ & $78{,}24\%$ & $85{,}03\%$ \\
        warm\_iter15 (vòng 15) & $59{,}40\%$ & $80{,}33\%$ & $78{,}57\%$ & $85{,}39\%$ \\
        warm\_iter20 (vòng 20) & $59{,}35\%$ & $80{,}69\%$ & $78{,}62\%$ & $85{,}70\%$ \\
        warm\_iter25 (vòng 25) & $59{,}44\%$ & $80{,}98\%$ & $77{,}84\%$ & $85{,}19\%$ \\
        warm\_iter30 (vòng 30) & $59{,}69\%$ & $80{,}69\%$ & $78{,}29\%$ & $85{,}37\%$ \\
        warm\_iter35 (vòng 35) & $59{,}33\%$ & $80{,}58\%$ & $78{,}02\%$ & $85{,}08\%$ \\
        warm\_iter40 (vòng 40) & $59{,}62\%$ & $80{,}96\%$ & $78{,}26\%$ & $85{,}90\%$ \\
        \midrule
        \multicolumn{5}{@{}l}{\textit{Nhánh Scratch (khởi tạo ngẫu nhiên):}} \\
        scratch\_iter00 (vòng 0) & $33{,}67\%$ & $71{,}71\%$ & $81{,}45\%$ & $81{,}48\%$ \\
        scratch\_iter05 (vòng 5) & $42{,}09\%$ & $74{,}76\%$ & $74{,}96\%$ & $80{,}32\%$ \\
        scratch\_iter10 (vòng 10) & $44{,}19\%$ & $76{,}21\%$ & $75{,}36\%$ & $82{,}21\%$ \\
        scratch\_iter15 (vòng 15) & $44{,}61\%$ & $76{,}58\%$ & $74{,}65\%$ & $81{,}52\%$ \\
        scratch\_iter20 (vòng 20) & $45{,}96\%$ & $77{,}21\%$ & $75{,}03\%$ & $81{,}77\%$ \\
        scratch\_iter25 (vòng 25) & $46{,}41\%$ & $77{,}18\%$ & $75{,}16\%$ & $81{,}70\%$ \\
        scratch\_iter30 (vòng 30) & $47{,}74\%$ & $77{,}59\%$ & $75{,}43\%$ & $82{,}25\%$ \\
        scratch\_iter35 (vòng 35) & $47{,}39\%$ & $77{,}92\%$ & $74{,}96\%$ & $81{,}94\%$ \\
        scratch\_iter40 (vòng 40) & $47{,}79\%$ & $77{,}99\%$ & $75{,}58\%$ & $82{,}48\%$ \\
        \bottomrule
    \end{tabular}
\end{table}

\section{Kết quả Giải đấu Vòng tròn Bản trước (Lịch sử)}
Bảng~\ref{tab:appendix_old_rr} lưu giữ kết quả giải đấu vòng tròn mốc lịch sử (seed 20261121, nguồn: \texttt{benchmarks/round\_robin/round\_robin\_final\_colab\_agents\_seed20261121\_v2.json}) sử dụng mô hình CNN v1 mốc cơ sở cũ và PUCT-AZ lịch sử để đối chứng quá trình phát triển.

\begin{table}[H]
    \centering
    \small
    \caption{Kết quả giải đấu vòng tròn bản trước (Mốc cơ sở seed 20261121, loại 6 khai cuộc).}
    \label{tab:appendix_old_rr}
    \begin{tabular}{@{} l r r r r r @{}}
        \toprule
        \textbf{Cặp đấu (T1 vs T2)} & \textbf{Thắng} & \textbf{Hòa} & \textbf{Thua} & \textbf{Score (T1)} & \textbf{CI 95\% Bootstrap} \\
        \midrule
        Random vs Heuristic & $0$ & $0$ & $100$ & $0{,}000$ & $[0{,}000;\; 0{,}000]$ \\
        Random vs AlphaBeta-4 & $0$ & $0$ & $100$ & $0{,}000$ & $[0{,}000;\; 0{,}000]$ \\
        Random vs MCTS-500 & $0$ & $0$ & $100$ & $0{,}000$ & $[0{,}000;\; 0{,}000]$ \\
        Heuristic vs AlphaBeta-4 & $22$ & $4$ & $74$ & $0{,}240$ & $[0{,}165;\; 0{,}315]$ \\
        Heuristic vs MCTS-500 & $31$ & $7$ & $62$ & $0{,}345$ & $[0{,}270;\; 0{,}420]$ \\
        Heuristic vs CNN-Policy (v1) & $49$ & $11$ & $40$ & $0{,}545$ & $[0{,}460;\; 0{,}635]$ \\
        AlphaBeta-4 vs MCTS-500 & $56$ & $6$ & $38$ & $0{,}590$ & $[0{,}510;\; 0{,}670]$ \\
        AlphaBeta-4 vs CNN-Policy (v1) & $67$ & $4$ & $29$ & $0{,}690$ & $[0{,}600;\; 0{,}775]$ \\
        AlphaBeta-4 vs PUCT-AZ (cũ) & $45$ & $9$ & $46$ & $0{,}495$ & $[0{,}405;\; 0{,}585]$ \\
        MCTS-500 vs CNN-Policy (v1) & $60$ & $6$ & $34$ & $0{,}630$ & $[0{,}540;\; 0{,}715]$ \\
        MCTS-500 vs PUCT-AZ (cũ) & $32$ & $7$ & $61$ & $0{,}355$ & $[0{,}275;\; 0{,}440]$ \\
        CNN-Policy (v1) vs PUCT-AZ (cũ) & $22$ & $3$ & $75$ & $0{,}235$ & $[0{,}165;\; 0{,}305]$ \\
        \bottomrule
    \end{tabular}
\end{table}

\section{Kết quả Ba kiểm tra Cài đặt Tự chơi (Sanity Checks)}
Thực hiện chạy kiểm tra tự động trên mã nguồn cục bộ:
\begin{enumerate}
    \item \textbf{Kiểm tra (a) --- Dấu giá trị ở thế thắng 1 nước và thua 2 nước:}
    Trạng thái kết thúc khi người chơi 1 thắng gán giá trị $+1{,}0$ theo góc nhìn người thắng. Trong \texttt{BatchedPUCT}, do trạng thái kết thúc bảo toàn người thắng là \texttt{current\_player}, nên khi lan truyền ngược (backpropagation) về nút cha, quan hệ \texttt{node.current\_player == parent.current\_player} được thỏa mãn; giá trị $+1{,}0$ được giữ nguyên dấu dương cho nút cha, giúp PUCT chọn chính xác nước đi chiến thắng. Ở thế thua 2 nước (chưa qua huấn luyện), PUCT thô không tự chặn do thiếu tìm kiếm sâu, nhưng bộ lọc an toàn ép buộc chặn thành công.
    \item \textbf{Kiểm tra (b) --- Tương đương suy luận đơn lẻ và theo lô:}
    So sánh đầu ra giữa suy luận từng trạng thái độc lập và suy luận theo lô (\textit{batched inference}) trên các trạng thái bàn cờ khác nhau: độ lệch cực đại đối với đầu ra chính sách là $1{,}86 \times 10^{-8}$ và đầu ra giá trị là $7{,}45 \times 10^{-9}$. Cả hai đầu ra đều hoàn toàn tương đương trong ngưỡng sai số số thực ($<10^{-5}$).
    \item \textbf{Kiểm tra (c) --- Tính nhất quán hướng bàn cờ giữa Engine Grid và Bitboard:}
    Đặc trưng 2 mặt phẳng sinh ra từ hàm \texttt{encode\_planes} trên trạng thái bàn cờ của \texttt{ConnectFour} và \texttt{BitboardConnectFour} hoàn toàn trùng khớp từng phần tử (\texttt{strictly equal}, kích thước $2 \times 6 \times 7$), xác nhận tính nhất quán hướng biểu diễn.
\end{enumerate}

\section{Truy xuất nguồn dữ liệu và hiện vật thực nghiệm}
Bảng~\ref{tab:artifact_traceability} ánh xạ các kết quả chính tới dữ liệu,
metadata và hiện vật được lưu trong repo. Tập kiểm thử đóng băng có tệp vị trí
và manifest riêng; các lần đánh giá tham chiếu cùng SHA256 và commit solver.

\small
\begin{longtable}{@{}p{3.4cm}p{7.1cm}p{4.0cm}@{}}
    \caption{Ánh xạ số liệu báo cáo với các hiện vật và nguồn kiểm chứng trong repo.}
    \label{tab:artifact_traceability}\\
    \toprule
    \textbf{Thành phần / Nội dung} & \textbf{Đường dẫn tệp nguồn trong repo} & \textbf{Mã định danh / Ghi chú} \\
    \midrule
    \endfirsthead

    \multicolumn{3}{l}{\textit{Bảng~\thetable{} (tiếp theo)}}\\
    \toprule
    \textbf{Thành phần / Nội dung} & \textbf{Đường dẫn tệp nguồn trong repo} & \textbf{Mã định danh / Ghi chú} \\
    \midrule
    \endhead

    \midrule
    \multicolumn{3}{r}{\textit{Tiếp tục ở trang sau}}\\
    \endfoot

    \bottomrule
    \endlastfoot

    Tập kiểm thử solver đóng băng
    & \texttt{results/solver\_eval\_seed20261122.positions.json};
      \texttt{results/solver\_eval\_seed20261122.manifest.json}
    & 1.500 vị trí; SHA256
      \texttt{3d8f9c747889d44821d8fd2ebe338570f59fef67f7e8961dd05153ef82d43b01};
      Pons commit
      \texttt{d6ba50d8aaf2308c769d9bf2abd42d90f34baf41} \\

    Đánh giá solver chính thức
    & \texttt{results/solver\_eval\_phase\_c\_local\_20261123.json};
      \texttt{results/solver\_eval\_cnn\_v3\_soft\_3seeds\_20261122.json}
    & Bảng solver chính và kết quả CNN v3 ba seed; bootstrap 10.000 lần theo cụm ván nguồn \\

    Mốc CNN v2 trước CNN v3
    & \texttt{results/solver\_eval\_cnn\_v2\_3seeds.json};
      \texttt{docs/LEARNING\_RESULTS.md}
    & Kết quả CNN v2 ba seed, dùng làm mốc lịch sử \\

    Lựa chọn và ablation CNN v3 trên validation
    & \texttt{results/cnn\_v3\_validation\_selection\_20261122.json};
      \texttt{docs/CNN\_V3\_SELECTION.md}
    & 5 cấu hình, 3 seed/cấu hình, 4.530 vị trí validation/seed; chọn Pons-soft 100\% \\

    Hiện vật huấn luyện CNN v3
    & \texttt{results/colab/cnn\_v3\_30k\_seed20261121/};
      \texttt{results/colab/cnn\_v3\_30k\_seed20261122/};
      \texttt{results/colab/cnn\_v3\_30k\_seed20261123/}
    & Manifest, train log và checkpoint tốt nhất được lưu cho các run; source bundle v11 không có trong repo. Manifest ghi source-code SHA256, không ghi commit \\

    Round-Robin mốc lịch sử
    & \texttt{docs/ROUND\_ROBIN\_FINAL.md};
      \texttt{benchmarks/round\_robin/round\_robin\_final\_colab\_agents\_seed20261121\_v2.json};
      \texttt{benchmarks/round\_robin/round\_robin\_final\_colab\_agents\_seed20261121\_v2.meta.json};
      \texttt{benchmarks/round\_robin/round\_robin\_final\_colab\_agents\_seed20261121\_v2.games.jsonl}
    & 8 tác tử, 100 ván/cặp, 50 khai cuộc 6 ply ghép cặp; fair filter giữ 50/56 khai cuộc thử \\

    Round-Robin Phase C
    & \texttt{docs/ROUND\_ROBIN\_PHASE\_C\_20261123.md};
      \texttt{results/round\_robin\_phase\_c\_20261123.json};
      \texttt{results/round\_robin\_phase\_c\_20261123.meta.json};
      \texttt{results/round\_robin\_phase\_c\_20261123.games.jsonl}
    & 8 tác tử, 28 cặp, 100 ván/cặp; 50 khai cuộc ghép cặp. Giữ 50/61 lần thử; loại 10 opening có thắng ngay và 1 trạng thái trùng \\

    Đường cong validation Phase C V2
    & \texttt{results/phase\_c\_local/validation\_curve\_20261123.json};
      \texttt{docs/PHASE\_C\_V2\_RESULTS.md};
      \texttt{docs/report/assets/phase\_c\_validation\_curve.png}
    & 40 vòng mỗi nhánh; 18 snapshot. Source-code SHA256
      \texttt{703fd34468d0038e328078495ebb67e55104c4cdfbfde3941d90ce988c5ef8} \\

    Manifest và checkpoint Phase C V2
    & \texttt{results/phase\_c\_local/az\_v2\_warm\_seed20261121/};
      \texttt{results/phase\_c\_local/az\_v2\_scratch\_seed20261122/}
    & Hai run, mỗi run 40 vòng; manifest, log huấn luyện và snapshot checkpoint \\

    Source snapshot Phase C
    & \texttt{results/phase\_c\_local/source\_snapshot\_703fd34468d0038e328078495ebb67e55104c4cdfbfde3941e4d90ce988c5ef8.zip};
      tệp manifest cùng tên
    & Archive SHA256
      \texttt{2aff20dc96fa1c9129c96e607d25816ef9ef0e33deab21417758208c7c331338};
      không có source bundle v12 \\

    Chẩn đoán lỗi CNN và value head
    & \texttt{results/cnn\_first\_errors\_seed20261122.json};
      \texttt{docs/CNN\_LOSS\_ANALYSIS.md};
      \texttt{results/value\_head/value\_cnn\_seed20261121\_20261121.json}
    & Chẩn đoán trên tập ván thua đã phân tích; value head được đánh giá riêng \\

    Hồ sơ tài nguyên và môi trường
    & \texttt{results/resource\_profile\_20261121\_v2.json};
      \texttt{docs/RESOURCE\_PROFILE.md};
      \texttt{docs/COLAB.md}
    & Tách số đo local khỏi thông tin tài nguyên Colab; một số manifest lịch sử không ghi peak VRAM \\
\end{longtable}
\normalsize

\section{Danh mục các mục cần hoàn thiện trước khi nộp}
Các vị trí giữ chỗ hành chính cần người dùng điền trước khi nộp:
\begin{enumerate}
    \item \textbf{Thông tin hành chính trên Trang bìa:} Điền tên Trường, Khoa, Họ tên Giảng viên hướng dẫn, Mã nhóm và liên kết kho lưu trữ GitHub tại các nhãn \texttt{[CẦN BỔ SUNG]}.
    \item \textbf{Danh sách thành viên và Tỷ lệ đóng góp:} Nếu thực hiện theo nhóm, bổ sung bảng phân công nhiệm vụ và tỷ lệ phần trăm đóng góp của các thành viên.
    \item \textbf{Hình ảnh minh họa Giao diện Demo:} Bổ sung ảnh chụp màn hình chế độ chơi Người--Máy và Máy--Máy nếu có yêu cầu trong rubric.
    \item \textbf{Kiểm tra thư viện tài liệu tham khảo:} Rà soát lại các mục đánh dấu \texttt{[CẦN KIỂM TRA]} trong danh mục tài liệu tham khảo.
\end{enumerate}

\end{document}
